% =======================================================================
\section{Results}

\subsection{Connected components (\texorpdfstring{$H_0$}{})}
\label{sec:scaleh0}

Addressing question~\ref{q:levels}, the uniform subsample reproduces
the three group-count levels of Section~\ref{sec:cutoffselection}.
At filtration zero, $H_0$ equals the group count $k$ by construction,
so the check is that the subsample matches the full-sweep counts.
Those levels are individual
($\delta = 2.75$~m), tactical ($\delta = 11.75$~m), and team
($\delta = 23.0$~m). Throughout, $H_0$ is the group count: the number
of centroids after clustering. Equivalently, it is $\beta_0$ of the
Vietoris--Rips complex on those centroids at filtration zero.
Primary-match figures come from the uniform 150-frame sample of match
1996435 (Section~\ref{sec:data}). Ten-match aggregates use the same
sampling scheme (1{,}500 frames in total).

At the individual level,
$k = 19.32 \pm 2.34$ (mean $\pm$ s.d.; range 10--22). At the tactical
level, $k = 5.04 \pm 1.69$ (range 2--10). At the
team level, the cloud usually collapses to one or
two groups: mean $1.96 \pm 0.36$. All 22 players lie in a single
group in 8.7\% of frames, split across two in 86.7\%, and across
three in 4.7\%.

The same three-level pattern holds in all ten matches. Grand means are
individual $k = 19.31 \pm 0.31$, tactical $k = 5.09 \pm 0.38$, and
team $k = 1.98 \pm 0.06$ (mean $\pm$ s.d.\ across matches).
Every match stays within the acceptance intervals of
Section~\ref{sec:cutoffselection}: the ten per-match means all fall
inside $[15,22]$, $[4,10]$, and $[1,2.5]$
(Figure~\ref{fig:cutoff}b). Figure~\ref{fig:cutoff} records the sweep
from which those cutoffs were taken. The pooled curve is asymmetric:
about $1.6$ groups merge per metre from the individual to the
tactical level, but only about $0.27$ per metre from the tactical to
the team level, so the two envelopes remain over a wide band
of $\delta$ (Figure~\ref{fig:cutoff}a). The spatial realisation on a
single frame (Figure~\ref{fig:cutoff}d--f) shows the same three
regimes: a mostly resolved set of agents at $2.75$~m, a handful of
coordinating groups at $11.75$~m, and two spatial envelopes at
$23.0$~m, which are topological components rather than the two teams.

\subsection{\texorpdfstring{$H_1$}{} loop detection}
\label{sec:h1detection}

Addressing question~\ref{q:homology}, scale-attributable homology
detects loop structure at two of the three recovered levels. On
the primary match, the pipeline detects 403 $H_1$ loops across 150
uniformly sampled frames (Table~\ref{tab:h1single}). Loops at the
individual level are frequent. They appear in 144 of 150 frames
(96.0\%), at a mean rate of 2.52 per frame, with short filtration
lifetime ($\bar{p}=1.931 \pm 1.735$~m). Loops at the tactical level are
rarer: 23 frames (15.3\%), mean rate 0.17 per frame, with longer
filtration lifetime ($\bar{p}=3.778 \pm 2.723$~m, maximum 10.771~m).
Loops at the team level never occur ($H_1 = 0$ across all 1{,}500
frames in the ten-match sample).

Table~\ref{tab:h1single} collects these statistics in compact form.
$F$ denotes presence: the fraction of frames with at least one loop;
L/f, the mean loop count per frame; $\bar{p}$, mean filtration
lifetime (death minus birth over loops, in metres); and $p_{\max}$,
maximum filtration lifetime.

\begin{table}[h]
\centering
\caption{Single-match \texorpdfstring{$H_1$}{} statistics (primary match, 150 frames).
$\bar{p}$ is the mean of death minus birth over loops, in metres on the
Vietoris--Rips parameter, not a duration in time.}
\label{tab:h1single}
\begin{tabular}{@{}lccccr@{}}
\toprule
Level & Total & $F$ (\%) & L/f & $\bar{p}$ (m) & $p_{\max}$ (m) \\
\midrule
Individual (2.75~m) & 378 & 144/150 (96.0\%) & 2.52 & $1.931\pm1.735$ & 12.991 \\
Tactical (11.75~m)  & 25  & 23/150 (15.3\%)  & 0.17 & $3.778\pm2.723$ & 10.771 \\
Team (23.0~m)       & 0   & 0/150 (0\%)      & N/A  & N/A             & N/A \\
\bottomrule
\end{tabular}
\end{table}

\begin{remark}[Team-level $H_1$ needs four centroids]
\label{rem:teamnull}
A Vietoris--Rips complex needs at least four vertices to carry a
non-trivial 1-cycle. Three pairwise-connected vertices span a
2-simplex as soon as the last edge enters, so $k\le 3$ cannot carry a
loop. At $\delta = 23.0$~m the 436{,}648 complete frames split as
$k\le 2$ in 93.6\%, $k=3$ in 6.4\%, and $k=4$ in the remaining 141
frames (0.032\%). No frame has $k>4$. A loop is therefore impossible
on all but those 141 frames. On the 141 frames the computed $H_1$
count is zero, so team-level $H_1$ is absent on the whole sweep.
The 1{,}500-frame sample contains one of those $k=4$ frames, and that
frame has no loop either. Team-level loop structure on the typical
envelope would require a different representation (density fields or
Delaunay triangulations, for example).
\end{remark}

The two-level $H_1$ pattern extends to all ten matches (1{,}500
uniformly sampled frames; Table~\ref{tab:h1multi}). Individual-level
presence is $96.5\%\pm1.5\%$ (95\% CI $95.6$--$97.3\%$), between 94.7\%
and 98.7\% in every match. Tactical presence is $18.8\%\pm6.5\%$ (95\%
CI $15.2$--$22.7\%$), ranging from 11.3\% to 31.3\% across matches. The
primary match (15.3\% tactical) sits within one standard deviation of
the ten-match mean. The match was chosen for broadcast quality and
event annotation, not for its tactical $H_1$ rate.
Conditional mean filtration lifetime on the ten-match sample is
$1.886$~m at the individual level and $3.411$~m at the tactical level,
the same ordering as Table~\ref{tab:h1single}.

\begin{table}[h]
\centering
\caption{Multi-match \texorpdfstring{$H_1$}{} statistics (10 matches).
Presence = mean \texorpdfstring{$\pm$}{} s.d.\ across matches; 95\%
bootstrap CIs from 1{,}000 resamples over matches.
$\bar{p}$ is mean filtration lifetime over loops
(\texorpdfstring{$\sum p / \sum H_1$}{}), not the per-frame average
that includes zeros. The s.d.\ column is the standard deviation of the
ten per-match conditional means.}
\label{tab:h1multi}
\begin{tabular}{@{}lccccr@{}}
\toprule
Level & Total & Presence (\%) & 95\% CI & $\bar{p}$ (m) & s.d. \\
\midrule
Individual & 4{,}039 & $96.5\%\pm1.5\%$ & $95.6$--$97.3\%$ & 1.886 & 0.145 \\
Tactical   & 306     & $18.8\%\pm6.5\%$ & $15.2$--$22.7\%$ & 3.411 & 0.419 \\
Team       & 0       & $0.0\%\pm0.0\%$  & N/A              & N/A   & N/A \\
\bottomrule
\end{tabular}
\end{table}

Presence alone cannot separate arrangement from group count, because
more centroids close more cycles. The matched null of
Section~\ref{sec:stats} holds both the group count and the spatial
envelope fixed and randomises only the arrangement. Against that null,
tactical presence is 18.8\% compared with 9.1\%, an excess of $+9.7$
percentage points (Table~\ref{tab:null}). Table~\ref{tab:nullk} shows
where the excess arises. No loops appear in the 40.6\% of frames with
four or fewer tactical groups. Presence then rises from 9.3\% at
$k=5$ to 67.0\% at $k=8$ and stays above the null at every such $k$.
At the individual level the null rate is already 91.4\%, since about
twenty centroids in a bounded region almost always close a cycle, so
the observed excess is modest ($+5.1$ percentage points).

\begin{table}[h]
\centering
\caption{\texorpdfstring{$H_1$}{} presence against a group-count- and envelope-matched null (10 matches, 1{,}500 frames, 200 null replicates per frame). Excess is in percentage points, with 95\% bootstrap CIs over matches.}
\label{tab:null}
\begin{tabular}{@{}lcccc@{}}
\toprule
Level & Observed & Null & Excess & 95\% CI \\
\midrule
Individual & 96.5\% & 91.4\% & $+5.1$ & $[+3.9, +6.1]$ \\
Tactical   & 18.8\% & 9.1\%  & $+9.7$ & $[+7.3, +12.2]$ \\
\bottomrule
\end{tabular}
\end{table}

\begin{table}[h]
\centering
\caption{Tactical-level \texorpdfstring{$H_1$}{} presence by group count \texorpdfstring{$k$}{}, against the matched null of Table~\ref{tab:null} (10 matches, 1{,}500 frames).}
\label{tab:nullk}
\begin{tabular}{@{}rcccc@{}}
\toprule
$k$ & Frames & Observed & Null & Excess \\
\midrule
$\le 4$ & 609 & $0.0\%$ & $0.4\%$ & $-0.4$ \\
5 & 344 & $9.3\%$ & $4.1\%$ & $+5.2$ \\
6 & 252 & $26.6\%$ & $12.0\%$ & $+14.5$ \\
7 & 149 & $51.0\%$ & $21.7\%$ & $+29.3$ \\
8 & 94 & $67.0\%$ & $33.4\%$ & $+33.6$ \\
$\ge 9$ & 52 & $84.6\%$ & $49.0\%$ & $+35.6$ \\
\bottomrule
\end{tabular}
\end{table}

\subsection{Closed cycle structures}
\label{sec:cyclegeometry}

All 403 primary-match $H_1$ bars have a simple-cycle representative
from the reduction in Section~\ref{sec:cycles}. The cycle is present
at birth by construction: its longest edge is the birth value. The cycles have four
to thirteen vertices, and 352 of the 403 have four, five, or six.
Figure~\ref{fig:cycles} shows sample frame~95 at both levels. The left
column is the vertex set after clustering. The right column is the
representative of largest $p$ on that set.
At $\delta=2.75$~m the cycle has eight vertices, seven of them
singleton clusters and one a pair, and the edges run from $6.9$ to
$15.6$~m (\texorpdfstring{$p=6.540$}{}~m).
At $\delta=11.75$~m the cycle has six groups, one a $14$-player
single-linkage cluster, and the edges run from $14.6$ to $22.1$~m
(\texorpdfstring{$p=9.942$}{}~m).
The two cycles are not the same bar: they use different vertices and
are born at $15.6$~m and $22.1$~m.
Figure~\ref{fig:h1bars} shows those bars on the persistence diagram,
together with every finite $H_1$ bar on the primary match.
The bars of largest $p$ in Table~\ref{tab:h1single} are not used for
the figure: the individual-level bar of largest $p$ (frame~141) is a
near-collinear sliver.

\begin{figure}[tp]
\centering
\includegraphics[width=\textwidth,keepaspectratio]{figures/fig3_cycle_geometry}
\caption{Cutoff versus filtration on one sample frame of SkillCorner
match 1996435 (sample frame~95). Left: vertex set after clustering at
\texorpdfstring{$\delta$}{}. Gold markers belong to this row's cycle
clusters; the hull in (c) is the 14-player single-linkage cluster;
dashed circles have radius \texorpdfstring{$\delta$}{} on the cycle
vertices only. Right: \texorpdfstring{$H_1$}{} representative of
largest \texorpdfstring{$p$}{} on those vertices
(\texorpdfstring{$p=6.540$}{}~m and
\texorpdfstring{$p=9.942$}{}~m). Cycle edges are filtration
distances (\texorpdfstring{$6.9$--$15.6$}{}~m in (b);
\texorpdfstring{$14.6$--$22.1$}{}~m in (d)). The longest edge is the
birth value. Panel (d) draws centroids, not the absorbed players.
Eight vertices in (b) and six in (d), not the bar of largest
\texorpdfstring{$p$}{} on the match.
A 10~m scale bar is shown on each panel. Uniform sample,
\texorpdfstring{$n=150$}{} frames, every 290th complete frame (the
same rule as the ten-match analysis).}
\label{fig:cycles}
\end{figure}

\begin{figure}[tp]
\centering
\includegraphics[width=\textwidth,keepaspectratio]{figures/fig4_h1_diagrams}
\caption{Finite \texorpdfstring{$H_1$}{} bars at two organisational
levels on the primary match (uniform 150-frame sample).
\textbf{(a)}~Individual and \textbf{(b)}~tactical persistence diagrams
for sample frame~95, the same frame as Figure~\ref{fig:cycles}.
The marked pair is the bar of largest \texorpdfstring{$p$}{} on that frame
(\texorpdfstring{$p=6.540$}{}~m and
\texorpdfstring{$p=9.942$}{}~m), whose geometric representative appears
in Figure~\ref{fig:cycles}.
\textbf{(c)}~Every finite bar on the primary match, one track per
organisational level (378 individual, 25 tactical).
The dashed ticks are Table~\ref{tab:h1single}'s $\bar{p}$: the mean of
these bar lengths.
The two levels are not superimposed: that would recreate the mixing of
Section~\ref{sec:scaleconflation}.
The individual-level bar of largest \texorpdfstring{$p$}{} in
Table~\ref{tab:h1single} (frame~141) is a near-collinear sliver and is
not the cycle shown in Figure~\ref{fig:cycles}.}
\label{fig:h1bars}
\end{figure}

\subsection{Sensitivity analysis}
\label{sec:sensitivity}

Detection is not confined to the adopted cutoff alone. Tactical $H_1$ counts fall monotonically
as the cutoff widens. On the
primary match's 150 frames, sweeping $\delta\in[6,16]$~m gives 275
loops (87.3\% of frames) at $\delta=6$~m and none at $\delta=16$~m.
The operative range is $[6,14]$~m (Table~\ref{tab:sensitivity}). The
adopted cutoff $\delta=11.75$~m lies at the conservative end of that
range (15.3\% frame presence).

\begin{table}[h]
\centering
\caption{Tactical-level cutoff sensitivity (150 frames, primary match).}
\label{tab:sensitivity}
\begin{tabular}{@{}lccc@{}}
\toprule
$\delta$ (m) & Total $H_1$ & Frame presence & Mean $k$ \\
\midrule
6  & 275 & 87.3\% & 13.4 \\
8  & 162 & 68.0\% & 9.8 \\
10 & 78  & 42.7\% & 7.0 \\
11.75 & 25 & 15.3\% & 5.0 \\
12 & 21  & 12.7\% & 4.8 \\
14 & 5   & 3.3\%  & 3.5 \\
16 & 0   & 0.0\%  & 2.8 \\
\bottomrule
\end{tabular}
\end{table}

\subsection{Complementarity of the two \texorpdfstring{$H_1$}{} levels}
\label{sec:complementarity}

Addressing question~\ref{q:distinct}, the two $H_1$ levels carry
distinct rather than redundant information.
Over 1{,}500 uniformly sampled frames, total individual and tactical
$H_1$ persistence correlate weakly (Spearman $\rho=0.234$, $p<0.001$;
95\% bootstrap CI $[0.181, 0.276]$). The same holds for loop counts
($\rho=0.205$, $p<0.001$). The finding does not depend on whether the
comparison uses total persistence or loop counts.

Co-occurrence exceeds chance (Fisher odds ratio 6.12, $p=0.002$;
bootstrap CI $[2.32, 16.33]$), and the association is one-sided
(Table~\ref{tab:fisher}). In 1{,}167 frames an individual loop appears
without a tactical partner. In only two frames the reverse occurs.
Because individual loops are present in almost every frame, joint
presence is expected once a tactical loop forms, and it does not imply
that the two levels record the same structure. The weak rank
correlation therefore remains the main evidence for complementarity.

\begin{table}[h]
\centering
\caption{Frame counts of \texorpdfstring{$H_1$}{} presence at the two levels (1{,}500 frames).}
\label{tab:fisher}
\begin{tabular}{@{}lcc@{}}
\toprule
 & Tactical $H_1$ present & Tactical $H_1$ absent \\
\midrule
Individual $H_1$ present & 280 & 1{,}167 \\
Individual $H_1$ absent  & 2   & 51 \\
\bottomrule
\end{tabular}
\end{table}

Bottleneck distance between the two levels'
diagrams has median 1.456~m and 95th percentile 3.556~m. The
95th percentile is close to the ten-match tactical mean bar length,
3.411~m (Table~\ref{tab:h1multi}). Landscape $L^2$ distance has median 5.477.
The levels differ by roughly as much as their features are large, not
by a small perturbation of one shared pattern.

\subsection{Event correlation}
\label{sec:eventcorr}

Addressing question~\ref{q:validity}, total persistence moves with
on-ball events. The comparison is a construct-validity check against
measurement noise. It uses 103{,}856 SkillCorner event--topology pairs
from the ten matches.

For each event, $\Delta$ is the change in mean total persistence from
the five 1~Hz frames before the event to the five after it
(Section~\ref{sec:stats}). At the individual level, an on-ball
engagement is followed by a drop. The mean change is
$\Delta=-0.253$~m ($n=8{,}927$, Wilcoxon signed-rank $p<10^{-13}$).
Build-up is followed by a rise. The mean change is
$\Delta=+0.700$~m ($n=618$, Wilcoxon signed-rank $p<10^{-7}$). At the tactical level the mean changes are
$\Delta=-0.056$~m for engagement ($n=8{,}927$, of which $3{,}738$ are
exact zeros; Wilcoxon signed-rank $p=0.013$ on the remainder) and
$\Delta=+0.339$~m for build-up ($n=618$, of which $153$ are exact
zeros; $p<10^{-5}$).

The event classes were specified in advance, and the $p$-values are
nominal. Multiple-testing control and a football reading of these
shifts belong to a companion manuscript.
